\newglossaryentry{closedworld}{
    name={Closed-World},
    description={Auswertungsszenario, in dem der Klassifikator nur Seiten aus einer bekannten, endlichen Menge überwachter Seiten unterscheiden muss.}
}

\newglossaryentry{openworld}{
    name={Open-World},
    description={Auswertungsszenario mit einer überwachten Seitenmenge und einer offenen Sammelklasse für alle übrigen Seiten.}
}

\newglossaryentry{circuitpadding}{
    name={Circuit Padding},
    description={In Tor integriertes Verteidigungsframework, das zusätzliche Zellen in einen Circuit einfügt, um Verkehrsmuster zu verschleiern.}
}

\newglossaryentry{tamaraw}{
    name={Tamaraw},
    description={Festrate-Padding-Verteidigung, die Pakete zu festen Zeitabständen sendet und damit Webseiten-spezifische Timing-Muster nivelliert.}
}

\newglossaryentry{monitorhost}{
    name={Monitor-Host},
    description={Simulierter Tor-Client, der innerhalb der Shadow-Simulation Wikipedia-Seiten abruft und dabei Verkehrsspuren für die Website-Fingerprinting-Klassifikation aufzeichnet.}
}

\newglossaryentry{zim}{
    name={ZIM},
    description={Komprimiertes Container-Format für offline nutzbare Wikipedia-Snapshots, wird in dieser Arbeit von \texttt{zimsrv.py} ausgeliefert.}
}

\newglossaryentry{shadow}{
    name={Shadow},
    description={Diskreter Netzwerkereignis-Simulator, der unmodifizierte Anwendungen wie den Tor-Daemon innerhalb einer deterministischen Simulation ausführt und in dieser Arbeit als Laborumgebung für die Website-Fingerprinting-Experimente dient.}
}

\newglossaryentry{klasse}{
    name={Klasse},
    description={Im Kontext dieser Arbeit eine zu unterscheidende Webseite. Closed-World mit 80 Klassen bedeutet 80 verschiedene Wikipedia-Artikel, die der Klassifikator unterscheiden muss.}
}
